% ==============================================================
% Formale Semantik des Ontophänomenalismus (FSO)
% Version 1.0
% Kompiliert mit: pdfLaTeX
% ==============================================================

\documentclass[11pt,a4paper]{article}

% ---------- Grundpakete ----------
\usepackage[utf8]{inputenc}
\usepackage[T1]{fontenc}
\usepackage[german]{babel}
\usepackage{amsmath,amssymb,amsthm}
\usepackage{geometry}
\usepackage{parskip}
\usepackage{enumitem}
\usepackage{booktabs}
\usepackage{xcolor}
\usepackage{hyperref}

% ---------- Layout ----------
\geometry{margin=2.5cm}
\setlist{itemsep=0.2em, topsep=0.2em}
\hypersetup{
    colorlinks=true,
    linkcolor=blue,
    urlcolor=blue,
    pdftitle={Formale Semantik des Ontoph\"{a}nomenalismus (FSO) -- Version 1.0},
    pdfauthor={Forschungsprogramm Ontoph\"{a}nomenalismus}
}

% ---------- Theorem-Umgebungen ----------
\theoremstyle{plain}
\newtheorem{theorem}{Theorem}[section]
\newtheorem{proposition}{Proposition}[section]
\newtheorem{corollary}{Korollar}[section]

\theoremstyle{definition}
\newtheorem{definition}{Definition}[section]
\newtheorem{remark}{Bemerkung}[section]
\newtheorem{methodennote}{Methodennote}[section]
\newtheorem{konvention}{Konvention}[section]

% ---------- Titel ----------
\title{
  \textbf{Formale Semantik des Ontoph\"{a}nomenalismus (FSO)}\\[0.4em]
  {\large Diskursdom\"{a}nen, Referenztheorie und sprachlogische Grundlagen}\\[0.5em]
  {\normalsize Version~1.0}
}
\author{
  \textbf{Forschungsprogramm Ontoph\"{a}nomenalismus}\\
  \small Siehe TRANSPARENCY f\"{u}r methodologische Urheberschaft und KI-Einsatz.
}
\date{Juli 2026}

\begin{document}

\maketitle

\begin{abstract}
Dieses Dokument legt die formalsemantischen Grundlagen des Ontoph\"{a}nomenalismus (OP)
fest. Es steht zwischen den methodologischen Konventionen (MKO) und den mathematischen
Grundlagen (MGO) und behandelt die Bedingungen, unter denen \"{u}ber ontologische
Gegenst\"{a}nde \"{u}berhaupt sinnvoll gesprochen werden kann.

FSO verbindet klassische Logik, Montague-Grammatik, Pr\"{a}suppositionstheorie,
Diskursdom\"{a}nen und dynamische Semantik zu einem gemeinsamen methodischen Rahmen.
Sein zentrales Ziel ist die Identifikation und Vermeidung von Kategorienfehlern ---
insbesondere der Reifikation negativer Grenzbegriffe wie dem \glqq{}absoluten
Nichts\grqq{}.

FSO dient allen nachfolgenden Dokumenten als sprachlogisches Fundament. Es erm\"{o}glicht
pr\"{a}zise Aussagen \"{u}ber Existenz, Referenz und Bestimmtheit, ohne die Fallen
nat\"{u}rlicher Sprache zu \"{u}bernehmen.

Es gelten die methodologischen Konventionen gem\"{a}\ss{} MKO.
\end{abstract}

\tableofcontents
\newpage

%==============================================================
\section{Zweck und Einordnung}
%==============================================================

\begin{methodennote}[Verh\"{a}ltnis zu MKO und MGO]
MKO definiert Sprachregeln und epistemische Standards f\"{u}r alle Dokumente des OP.
FSO vertieft diese Regeln auf der Ebene formaler Semantik: Es kl\"{a}rt, was Begriffe
wie \glqq{}Existenz\grqq{}, \glqq{}Referenz\grqq{} und \glqq{}Bestimmtheit\grqq{} im
Rahmen des OP formal bedeuten. MGO baut auf FSO auf und formalisiert die
mathematischen Strukturen, die FSO sprachlogisch vorbereitet.

FSO ist kein linguistisches Dokument. Es ist ein philosophisch-logisches Werkzeug
zur Vermeidung von Kategorienfehlern --- insbesondere dort, wo nat\"{u}rliche Sprache
ontologische Fragen in die Irre f\"{u}hren kann.
\end{methodennote}

%==============================================================
\section{Grundbegriffe der formalen Semantik}
%==============================================================

\subsection{Assertion und Pr\"{a}supposition}

\begin{definition}[Assertion]
Eine \textbf{Assertion} ist das, was ein Satz ausdr\"{u}cklich behauptet --- sein
propositionaler Gehalt, der einen Wahrheitswert annehmen kann.
\end{definition}

\begin{definition}[Pr\"{a}supposition]
Eine \textbf{Pr\"{a}supposition} ist eine Voraussetzung, die erf\"{u}llt sein muss,
damit ein Satz \"{u}berhaupt sinnvoll einen Wahrheitswert annehmen kann. Pr\"{a}suppositionen
bleiben unter Negation erhalten.
\end{definition}

\begin{remark}[Pr\"{a}suppositionstest]
Der klassische Test: Ein Satz $S$ pr\"{a}supponiert $P$, wenn sowohl $S$ als auch
$\neg S$ die G\"{u}ltigkeit von $P$ voraussetzen.

\medskip
\textit{Beispiel:} \glqq{}Dieses Glas ist nicht leer.\grqq{}\\
Assertion: Das Glas ist nicht leer.\\
Pr\"{a}supposition: Es gibt ein bestimmtes Glas. --- Diese bleibt auch unter Negation
erhalten.

\medskip
\textit{Konsequenz f\"{u}r den OP:} Die Frage \glqq{}Warum ist etwas und nicht
nichts?\grqq{} pr\"{a}supponiert, dass \glqq{}nichts\grqq{} einen koh\"{a}renten
Alternativzustand bezeichnet. Genau diese Pr\"{a}supposition ist Gegenstand der
Debatte --- sie kann nicht stillschweigend in die Fragestellung eingebaut werden.
\end{remark}

\subsection{Referenz und Reifikation}

\begin{definition}[Referenz]
\textbf{Referenz} bezeichnet die Beziehung zwischen einem sprachlichen Ausdruck und
dem Gegenstand, auf den er verweist. Erfolgreiche Referenz setzt voraus:
\begin{enumerate}
  \item Es gibt einen Referenten (ein Diskursobjekt).
  \item Identit\"{a}tsbedingungen f\"{u}r diesen Referenten sind definiert.
  \item \"{U}ber diesen Referenten kann sinnvoll pr\"{a}diziert werden.
\end{enumerate}
\end{definition}

\begin{definition}[Reifikation]
\textbf{Reifikation} bezeichnet den Kategorienfehler, einen Nicht-Gegenstand
(z.\,B.\ einen Grenzbegriff, ein Indefinitpronomen oder eine Quantifikation) wie
einen referenzf\"{a}higen Gegenstand zu behandeln.
\end{definition}

\begin{remark}[Reifikation des Nichts]
Die Substantivierung \glqq{}das Nichts\grqq{} ist eine klassische Reifikation. Durch
grammatische Substantivierung erh\"{a}lt \glqq{}nichts\grqq{} --- ein Indefinitpronomen
--- die scheinbare Form eines referenzf\"{a}higen Gegenstandes. Damit wird ihm bereits
genau jene minimale Bestimmtheit zugeschrieben, die sein Inhalt ausschlie\ss{}en soll.

Der korrekte Ausdruck ist: \glqq{}es existiert nichts\grqq{} --- als Quantifikation,
nicht als Referenz auf ein Objekt namens \glqq{}Nichts\grqq{}.
\end{remark}

%==============================================================
\section{Montague-Grammatik als formalsemantischer Rahmen}
%==============================================================

\begin{methodennote}[Status dieser Sektion --- MKO Ebene~3 und~4]
Die Montague-Grammatik wird hier als formales Werkzeug eingef\"{u}hrt. Ihre Anwendung
auf den OP ist ein konzeptuelles Modell im Sinne von MKO --- kein abgeschlossener
Beweis, sondern ein Pr\"{a}zisierungsrahmen.
\end{methodennote}

\begin{definition}[Montague-Semantik --- Grundprinzipien]
Die \textbf{Montague-Grammatik} (Richard Montague, 1970er) behandelt nat\"{u}rliche
Sprache mit denselben formalen Mitteln wie Logik. Ihre Kernprinzipien:
\begin{enumerate}
  \item \textbf{Typentheorie:} Jeder Ausdruck hat einen semantischen Typ.
        Individuen haben Typ $e$, Wahrheitswerte Typ $t$, Eigenschaften Typ
        $\langle e, t \rangle$, Quantoren Typ $\langle\langle e,t\rangle, t\rangle$.
  \item \textbf{Kompositionalit\"{a}t:} Die Bedeutung eines komplexen Ausdrucks
        ergibt sich regelm\"{a}\ss{}ig aus den Bedeutungen seiner Teile.
  \item \textbf{Intensionale Logik:} Bedeutungen sind nicht nur extensional
        (Mengen von Individuen), sondern intensional --- sie variieren \"{u}ber
        m\"{o}gliche Welten.
  \item \textbf{Generalisierte Quantoren:} Quantorausdr\"{u}cke wie \glqq{}nichts\grqq{},
        \glqq{}alles\grqq{}, \glqq{}einiges\grqq{} werden als Mengen von Eigenschaften
        behandelt, nicht als Referenzausdr\"{u}cke.
\end{enumerate}
\end{definition}

\begin{remark}[\glqq{}Nichts\grqq{} als generalisierter Quantor]
In der Montague-Semantik ist \glqq{}nichts\grqq{} ein generalisierter Quantor ---
kein Referenzausdruck. \glqq{}Nichts existiert\grqq{} bedeutet formal:
\[
  \neg \exists x\, E(x)
\]
wobei $E$ das Pr\"{a}dikat \glqq{}existiert\grqq{} bezeichnet. Diese Aussage
referiert auf kein Objekt namens \glqq{}Nichts\grqq{} --- sie quantifiziert
\"{u}ber eine Domäne und sagt, dass kein Element dieser Dom\"{a}ne die Eigenschaft
$E$ hat.

\textbf{Die entscheidende Konsequenz:} Der Existenzquantor $\exists$ ist niemals
dom\"{a}nenneutral. Er operiert stets \"{u}ber eine nicht-leere Diskursdom\"{a}ne.
Die Aussage $\neg \exists x\, E(x)$ ist formal undefiniert, wenn die Dom\"{a}ne
selbst leer oder inexistent ist.
\end{remark}

\begin{remark}[Definite Kennzeichnungen und Referenzanspruch]
Definite Nominalphrasen wie \glqq{}das absolute Nichts\grqq{} tragen in der
Montague-Semantik eine Pr\"{a}supposition: Es gibt genau einen Referenten der
entsprechenden Beschreibung. F\"{u}r \glqq{}das absolute Nichts\grqq{} sind die
Erf\"{u}llungsbedingungen dieser Pr\"{a}supposition prinzipiell nicht realisierbar
--- jeder Versuch, einen Referenten zu identifizieren, schreibt ihm bereits
Bestimmtheit zu.
\end{remark}

%==============================================================
\section{Hierarchie der Diskursdom\"{a}nen}
%==============================================================

\begin{definition}[Diskursdom\"{a}nen-Hierarchie]
Existenz, Bestimmtheit und Referenz sind stets relativ zu einer Diskursdom\"{a}ne
definiert. Die Dom\"{a}nen bilden eine Hierarchie zunehmender Einschr\"{a}nkung:
\[
  \mathcal{D}_{\mathrm{fakt}}
  \subset
  \mathcal{D}_{\mathrm{phys}}
  \subset
  \mathcal{D}_{\mathrm{log}}
  \subset
  \mathcal{D}_{\mathrm{kogn}}
  \subset
  \mathcal{D}_{\mathrm{synt}}
\]
\begin{itemize}
  \item $\mathcal{D}_{\mathrm{synt}}$: Alles syntaktisch Formulierbare.
  \item $\mathcal{D}_{\mathrm{kogn}}$: Alles begrifflich Repr\"{a}sentierbare.
  \item $\mathcal{D}_{\mathrm{log}}$: Alles logisch Konsistente.
  \item $\mathcal{D}_{\mathrm{phys}}$: Alles physikalisch M\"{o}gliche.
  \item $\mathcal{D}_{\mathrm{fakt}}$: Alles faktisch Existierende.
\end{itemize}
\end{definition}

\begin{proposition}[Implikationsrichtungen]
Aus der Hierarchie folgen einseitige Implikationen:
\begin{itemize}
  \item Was faktisch existiert, ist physikalisch m\"{o}glich.
  \item Was physikalisch m\"{o}glich ist, ist logisch konsistent.
  \item Was logisch konsistent ist, ist begrifflich repr\"{a}sentierbar.
  \item Was begrifflich repr\"{a}sentierbar ist, ist syntaktisch formulierbar.
\end{itemize}
Die Umkehrungen gelten im Allgemeinen \emph{nicht}: Syntaktische Formulierbarkeit
impliziert keine logische Konsistenz, logische Konsistenz keine physikalische
M\"{o}glichkeit, physikalische M\"{o}glichkeit keine faktische Existenz.
\end{proposition}

\begin{remark}[Praktische Konsequenz]
Diese Hierarchie ist das wichtigste Werkzeug zur Vermeidung von Kategorienfehlern
im OP. Wer \glqq{}das absolute Nichts\grqq{} syntaktisch formuliert, bewegt sich in
$\mathcal{D}_{\mathrm{synt}}$ --- nicht in $\mathcal{D}_{\mathrm{log}}$ oder
$\mathcal{D}_{\mathrm{phys}}$. Der Sprung von syntaktischer Formulierbarkeit zu
ontologischer M\"{o}glichkeit ist ein Kategorienfehler.
\end{remark}

%==============================================================
\section{Negative Dom\"{a}nen und ihre Grenzen}
%==============================================================

\begin{definition}[Negative Dom\"{a}ne]
Das Komplement der syntaktischen Dom\"{a}ne:
\[
  \mathcal{D}_{\mathrm{neg}} = \overline{\mathcal{D}_{\mathrm{synt}}}
\]
enth\"{a}lt alle Entit\"{a}ten $x$ f\"{u}r die gilt:
$\neg\,\mathrm{Symbolisierbar}(x)$.
\end{definition}

\begin{proposition}[Strukturlosigkeit negativer Dom\"{a}nen]
F\"{u}r jedes $x \in \mathcal{D}_{\mathrm{neg}}$ fehlen s\"{a}mtliche Voraussetzungen
f\"{u}r Referenz, Identit\"{a}tsbildung, Quantifikation und Pr\"{a}dikation. Negative
Dom\"{a}nen besitzen keine interne Struktur \"{u}ber die formallogisch sinnvoll
gesprochen werden k\"{o}nnte.
\end{proposition}

\begin{remark}[Das absolute Nichts als Grenzbegriff]
Das absolute Nichts ist kein Element irgendeiner positiven Dom\"{a}ne --- auch nicht
von $\mathcal{D}_{\mathrm{synt}}$. Sobald es syntaktisch formuliert wird, ist es
bereits in $\mathcal{D}_{\mathrm{synt}}$ eingetreten und hat damit genau die
Eigenschaft verloren, die es ausdr\"{u}cken soll.

Das absolute Nichts ist der Grenzbegriff an dem jede Dom\"{a}ne endet --- kein Punkt
innerhalb irgendeiner Dom\"{a}ne. Es bezeichnet den \"{U}bergang, nicht einen Zustand.
\end{remark}

%==============================================================
\section{Erweiterungen: Dynamische Semantik und Situation Semantics}
%==============================================================

\begin{methodennote}[Status dieser Sektion --- MKO Ebene~4 und~5]
Die folgenden Erweiterungen sind Analogien und Hypothesen im Sinne von MKO. Sie
skizzieren Forschungsrichtungen f\"{u}r die PST und zuk\"{u}nftige Versionen des FSO.
\end{methodennote}

\subsection{Dynamische Semantik}

\begin{remark}[Dynamische Semantik und Kontextaufbau --- MKO Ebene~4]
Die \textbf{Dynamische Semantik} (Groenendijk, Stokhof, Heim) behandelt Bedeutung
nicht als statische Relation zwischen Ausdruck und Referent, sondern als
\emph{Kontexttransformation}: Jeder Satz ver\"{a}ndert den Interpretationskontext
f\"{u}r nachfolgende S\"{a}tze.

F\"{u}r den OP ist dies relevant weil es zeigt: Diskursdom\"{a}nen entstehen durch
Diskurs --- sie sind nicht vorausgesetzt, sondern werden aufgebaut. Das korrespondiert
mit der Idee der Projektiven Reduktion (DPR): Metriken und Strukturen entstehen durch
Projektion, nicht als vorgegebene Rahmen.

\textbf{Hypothese (MKO Ebene~5):} Falls der OP zutrifft, k\"{o}nnte die dynamische
Semantik ein formales Modell f\"{u}r die Entstehung von Diskursdom\"{a}nen durch
Projektion liefern --- als Br\"{u}cke zwischen OPP und der raumzeitlichen Ebene.
\end{remark}

\subsection{Situation Semantics}

\begin{remark}[Situation Semantics und partielle Informationsstrukturen --- MKO Ebene~4]
Die \textbf{Situation Semantics} (Barwise, Perry) behandelt Bedeutung nicht relativ
zu vollst\"{a}ndigen m\"{o}glichen Welten, sondern relativ zu \emph{Situationen} ---
partiellen Informationsstrukturen, die nur einen Ausschnitt der Realit\"{a}t erfassen.

Dies korrespondiert direkt mit dem Fragedifferential (DFD): Eine individuierte
Perspektive $Q_E$ ist eine partielle Situation --- ein Ausschnitt des Gesamtfeldes $P$.
Das Fragedifferential $\mathcal{D}(E) = D_{\mathrm{KL}}(Q_E \| P)$ misst formal die
Distanz zwischen der partiellen Situation und der vollst\"{a}ndigen
Informationsstruktur.

\textbf{Hypothese (MKO Ebene~5):} Situation Semantics k\"{o}nnte ein formaler Rahmen
f\"{u}r die Theorie des Fragedifferentials sein: Individuation als Partialit\"{a}t einer
Situation, Erleben als relationale Struktur zwischen Situationen.
\end{remark}

\subsection{Typentheoretische Semantik}

\begin{remark}[Typentheoretische Semantik --- MKO Ebene~4]
Die \textbf{Typentheoretische Semantik} (Ranta, Martin-L\"{o}f) verbindet
Montague-Grammatik mit konstruktiver Mathematik. Typen entsprechen Propositionen,
Terme entsprechen Beweisen. Dies erm\"{o}glicht:
\begin{itemize}
  \item Pr\"{a}zise Typisierung aller Ausdr\"{u}cke des OP.
  \item Formale \"{U}berpr\"{u}fbarkeit von Schlussfolgerungen.
  \item Verbindung zur kategorientheoretischen Formulierung in MGO.
\end{itemize}
F\"{u}r den OP ist besonders relevant: In der konstruktiven Typentheorie existiert
ein Objekt genau dann, wenn ein Beweis (Term) f\"{u}r seinen Typ konstruiert werden
kann. Das korrespondiert mit der Unterscheidbarkeits-Pr\"{a}misse (UP): Existenz
erfordert konstruktive Bestimmtheit.
\end{remark}

\subsection{Freie Logik}

\begin{remark}[Freie Logik und Quantoren ohne Existenzpr\"{a}supposition --- MKO Ebene~4]
Die \textbf{Freie Logik} (Lambert, van Fraassen) erlaubt Quantoren \"{u}ber leere
oder nicht-existierende Terme. In klassischer Pr\"{a}dikatenlogik gilt:
$\exists x\, (x = a)$ f\"{u}r jeden Term $a$ --- jeder Term referiert auf etwas
Existierendes.

Freie Logik hebt diese Pr\"{a}supposition auf: Terme k\"{o}nnen auf nicht-existierende
Gegenst\"{a}nde referieren. Das erm\"{o}glicht formale Aussagen \"{u}ber
\glqq{}Nichts\grqq{} ohne Reifikation.

F\"{u}r den OP ist Freie Logik ein Kandidat f\"{u}r die pr\"{a}zise Behandlung von
Grenzbegriffen --- dort wo klassische Pr\"{a}dikatenlogik in die Reifikationsfalle
f\"{u}hrt. Ob Freie Logik die richtige Wahl ist, bleibt eine offene Forschungsfrage.
\end{remark}

%==============================================================
\section{Folgerung: Die Leibniz-Frage als kategoriale Fehlfrage}
%==============================================================

\begin{theorem}[Kategoriale Fehlfrage]
Die klassische Frage
\begin{quote}
\glqq{}Warum existiert \"{u}berhaupt etwas und nicht vielmehr nichts?\grqq{}
\end{quote}
ist in ihrer klassischen Form keine wohldefinierte ontologische Frage, sondern
enth\"{a}lt einen Kategorienfehler.
\end{theorem}

\begin{proof}[Konzeptuelle Herleitung]
\begin{enumerate}
  \item Die Frage setzt voraus, dass \glqq{}nichts\grqq{} einen koh\"{a}renten
        alternativen Gesamtzustand der Realit\"{a}t bezeichnet.
  \item Nach der Dom\"{a}nenhierarchie: \glqq{}Nichts\grqq{} als Gesamtzustand
        liegt au\ss{}erhalb jeder positiven Dom\"{a}ne einschlie\ss{}lich
        $\mathcal{D}_{\mathrm{synt}}$.
  \item Nach Montague: $\neg \exists x\, E(x)$ ben\"{o}tigt eine Dom\"{a}ne \"{u}ber
        die quantifiziert wird. Die totale Negation aller Existenz eliminiert simultan
        die Dom\"{a}ne --- der Ausdruck wird formal undefiniert.
  \item Nach der Pr\"{a}suppositionsanalyse: Die Frage pr\"{a}supponiert die
        G\"{u}ltigkeit von \glqq{}nichts\grqq{} als Alternative. Diese Pr\"{a}supposition
        kann nicht ohne Zirkel begr\"{u}ndet werden.
  \item Folglich: Die Frage kontrastiert Existenz mit einer Gr\"{o}\ss{}e, die die
        Voraussetzungen jeder Kontrastbildung selbst negiert.\qed
\end{enumerate}
\end{proof}

\begin{remark}[Was der OP stattdessen sagt]
Der Ontoph\"{a}nomenalismus beantwortet die Leibniz-Frage nicht. Er zeigt, warum sie
in ihrer klassischen Form logisch nicht wohldefiniert ist. Die pr\"{a}zisere Frage ---
die der OP stellt --- lautet:

\begin{quote}
\glqq{}Warum ist strukturierte Existenz notwendig --- und nicht eine
strukturlose Alternative?\grqq{}
\end{quote}

Diese Frage ist wohldefiniert: Sie fragt nicht nach \glqq{}Nichts\grqq{} als
Alternativzustand, sondern nach den Bedingungen der M\"{o}glichkeit von Struktur
\"{u}berhaupt. Die Antwort liefert die Kombination aus Axiom~1, Axiom~2 und der
Unterscheidbarkeits-Pr\"{a}misse (UP) in OP.
\end{remark}

\begin{remark}[Offene Forschungsfragen f\"{u}r die PST]
\begin{enumerate}
  \item Kann Freie Logik die Behandlung von Grenzbegriffen im OP pr\"{a}ziser
        formalisieren als klassische Pr\"{a}dikatenlogik?
  \item Wie verbindet sich Situation Semantics formal mit dem Fragedifferential
        (DFD)? Ist eine Perspektive formal eine Situation im Sinne von Barwise/Perry?
  \item Kann die Dom\"{a}nen-Hierarchie durch den Projektionsoperator (DPR) formal
        abgeleitet werden --- also zeigen dass Dom\"{a}nen entstehen statt vorausgesetzt
        werden?
  \item Welche Typen im Sinne der Martin-L\"{o}f-Typentheorie entsprechen den
        Grundbegriffen des OP (Attraktor, Projektion, Fragedifferential)?
\end{enumerate}
\end{remark}

%==============================================================
\section{Schnittstellen}
%==============================================================

\begin{itemize}
  \item \textbf{MKO (Methodologische Konventionen):} FSO konkretisiert die
        Sprachregeln von MKO auf der Ebene formaler Semantik. Insbesondere die
        Konvention gegen Reifikation und die f\"{u}nf Sprachebenen werden hier
        formal untermauert.
  \item \textbf{MGO (Mathematische Grundlagen):} FSO bereitet den
        kategorientheoretischen Rahmen von MGO vor. Typentheorie und
        Diskursdom\"{a}nen verbinden sich mit den topologischen und algebraischen
        Strukturen in MGO.
  \item \textbf{OP (Logische Grundlegung):} FSO liefert die formalsemantische
        Analyse von Axiom~1 (Notwendigkeit der Existenz), Axiom~2 (Exhaustivit\"{a}t)
        und der UP. Die Pr\"{a}suppositionsanalyse der Leibniz-Frage geh\"{o}rt zur
        Verteidigung von Axiom~1.
  \item \textbf{DFD (Das Fragedifferential):} Situation Semantics liefert einen
        m\"{o}glichen formalen Rahmen f\"{u}r das Fragedifferential als Ma\ss{} der
        Partialit\"{a}t einer Perspektive.
  \item \textbf{PST (Projektive Strukturtheorie):} Die offenen Forschungsfragen
        in Sektion~6 definieren ein Arbeitsprogramm f\"{u}r die PST im Bereich
        formaler Semantik.
\end{itemize}

\end{document}

